%% kapa: exam=1 round=21 year=2010 period=1 form=A subject=economics subject_ko=경제원론 seq=13 number=53
%% source: Q-Net 감정평가사 기출문제 (cst00309) artlSeq=5024181 file=extracted/1008/5024181/2010년 제21회 감정평가사 1차시험 기출문제/1교시.hwp
%% conversion: kordoc 4.24.51 -> tools/gen1.py (auto); review=proofread-sonnet; flags=none
\begin{question}{53}
수확체감의 법칙이 적용되는 성장이론에 관한 설명으로 옳지 않은 것은?
\begin{choices}
  \item 저축률이 증가하면 자원이 자본재 생산에 더 많이 투입된다.
  \item 저축률이 증가하면 일시적으로만 높은 성장률이 유지된다.
  \item 저축률이 증가하면 노동생산성이 높아지지만 장기적으로 노동생산성의 증가율은 높아지지 않는다.
  \item 다른 조건이 같을 경우 상대적으로 가난한 상태에서 출발하는 나라가 빠른 속도로 성장하기 어렵다.
  \item 국내저축뿐만 아니라 해외저축에 의한 투자도 생산증대와 실질임금 상승에 기여할 수 있다.
\end{choices}
\end{question}
\answer{4}
